\section{A square-root obstruction to uniform resolvent stability}\label{sec:conditioning}

The difficulty in \eqref{eq:ratio-error} is not removed merely by
imposing common source and zero-count bounds. The following family
makes the loss quantitative while fixing the translation scale.

Let \(H=H_0\), \(\alpha=(\log2)/2\), and
\[
 S_\alpha(u)=\frac{\Phi(u+\alpha)+\Phi(u-\alpha)}2,\qquad
 \Psi_\eta(u)=\frac{(1+\eta)\Phi(u)+(1-\eta)S_\alpha(u)}2,
\]
where \(0\le\eta\le\eta_0:=\tanh^2(\alpha/4)\).
The transforms are
\begin{equation}\label{eq:control}
 f_\eta(z)=\frac12\int_\R\Psi_\eta(u)\ee^{\ii zu}\dd u
 =H(z)M_\eta(z),\qquad
 M_\eta(z)=\cos^2(\alpha z/2)+\eta\sin^2(\alpha z/2).
\end{equation}
Every source is strictly positive, smooth and even, and
\(f_\eta(0)=H(0)\). The parameter \(\eta\) varies the source; it is
not Newman time. Each source has its own genuine heat orbit, but
\eqref{eq:control} is the time-zero factorization.

\begin{theorem}[Uniform square-root conditioning]\label{thm:conditioning}
The family \(\Psi_\eta\) has common bounds in every fixed exponentially
weighted source seminorm and common full zero-count constants. Its
relative source distance from \(\Psi_0\) is at most \(\eta\), and
each fixed safe-line logarithmic-resolvent jet changes by \(O(\eta)\),
uniformly along that line. Nevertheless,
\begin{equation}\label{eq:delta-scale}
 c\sqrt\eta\le\calJ[f_\eta](X)-\calJ[f_0](X)
                  \le C_\alpha\sqrt\eta
\end{equation}
for every \(X\in\R\), with constants independent of \(\eta,X\).
Thus no nondecreasing modulus \(o(\sqrt t)\), in particular no
uniform Lipschitz modulus, controls the Green difference from these
observations on a class containing this family.
\end{theorem}

\begin{proof}
The exact source difference is
\[
 \Psi_\eta-\Psi_0=\frac\eta2(\Phi-S_\alpha),\qquad
 |\Psi_\eta-\Psi_0|\le\eta\Psi_0.
\]
For fixed \(L,b\ge0\) and integers \(j,k\ge0\), its differentiated
weighted \(L^1\) norm is bounded by
\[
 \frac\eta2\int_\R |u|^j\ee^{Lu^2+b|u|}
            (|\Phi^{(k)}|+|S_\alpha^{(k)}|)\dd u.
\]
This integral is finite by Lemma~\ref{lem:source} and fixed translation.
The bound is independent of horizontal observation height.

Write
\[
 \ell=\frac{2\pi}{\alpha},\qquad
 x_k=\frac{(2k+1)\pi}{\alpha},\qquad
 d(\eta)=\frac2\alpha\atanh\sqrt\eta\le\frac12.
\]
At \(\eta=0\), the multiplier has double real zeros \(x_k\).
At \(\eta>0\), they are the simple zeros \(x_k\pm\ii d(\eta)\).
Any coincidence with an \(H\)-zero adds multiplicities. The fixed
lattice spacing gives a common unit-count bound and hence a common
full count after adding the canonical zeros.

The finite-factor argument of Section~\ref{sec:green}, applied
periodically to \(M_\eta\), gives
\begin{equation}\label{eq:lattice}
 \Delta_\eta(X):=\calJ[f_\eta](X)-\calJ[f_0](X)
       =\sum_{k\in\mathbb Z}G(X-x_k,d(\eta)),\qquad
 \calJ[f_0]=\calJ[H].
\end{equation}
One can justify the periodic limit directly on the cylinder: after
subtracting the Green singularities, the remaining function is
harmonic, has the correct boundary values and is bounded. Uniqueness
fixes its harmonic part. Alternatively \eqref{eq:G-decay} gives
absolute convergence of the periodized kernels. The real doubles of
\(M_0\) have zero interior contribution.

The nearest lattice point is at distance at most \(\ell/2\).
Using \eqref{eq:margin} for that point and
\eqref{eq:G-decay} for the complete lattice gives
\begin{equation}\label{eq:lattice-bounds}
 \frac{\sin a}{\cosh(a\ell/2)}\,d
 \le\Delta_\eta(X)
 \le K_G\coth(a\ell/2)\,d .
\end{equation}
The upper bound uses
\(\sup_X\sum_k\ee^{-a|X-x_k|}=\coth(a\ell/2)\).
Finally
\[
 \frac2\alpha\sqrt\eta\le d(\eta)
 \le\frac{2}{\alpha(1-\eta_0)}\sqrt\eta,
\]
which proves the two-sided Green scale.

It remains to verify the resolvent observations. With
\(R_\eta=f_\eta'/f_\eta\), \(R_*=f_0'/f_0\), and
\(B_\eta=(1+\eta)/(1-\eta)=\cosh(\alpha d)\), exact subtraction yields
\begin{equation}\label{eq:resolvent-difference}
 R_\eta-R_*=
 \frac{\alpha(B_\eta-1)\sin(\alpha z)}
 {(B_\eta+\cos(\alpha z))(1+\cos(\alpha z))}.
\end{equation}
The canonical resolvent cancels, but the multiplier's zeros remain
fully represented in \eqref{eq:lattice}. For \(y>1/2\) and
\(1\le B\le\cosh(\alpha/2)\),
\[
 |B+\cos(\alpha(x+\ii y))|\ge\cosh(\alpha y)-B,\qquad
 |\sin(\alpha(x+\ii y))|\le\cosh(\alpha y).
\]
For the first inequality, set \(t=\cos(\alpha x)\); the squared
modulus is \(t^2+2B\cosh(\alpha y)t+B^2+\sinh^2(\alpha y)\),
minimized at \(t=-1\).

Fix \(b>1\) and \(0<\varrho<b-1\), and define
\[
 C_{b,\varrho}=
 \frac{2\alpha\cosh(\alpha(b+\varrho))}
 {(1-\eta_0)
 [\cosh(\alpha(b-\varrho))-\cosh(\alpha/2)]
 [\cosh(\alpha(b-\varrho))-1]}.
\]
On every radius-\(\varrho\) disk centered at \(x+\ii b\),
\eqref{eq:resolvent-difference} is bounded by \(C_{b,\varrho}\eta\).
The actual resolvents are defined there by the zero strip.
Cauchy's estimate gives
\begin{equation}\label{eq:jet-change}
 \sup_x|\partial_z^j(R_\eta-R_*)(x+\ii b)|
 \le j!\varrho^{-j}C_{b,\varrho}\eta.
\end{equation}
Any fixed sum, maximum, or ordinary finite-dimensional norm of the
relative source distance and finitely many such jet seminorms is
therefore at most \(C_D\eta\), with weights independent of \(\eta\).
A nondecreasing \(\omega(t)=o(\sqrt t)\) would imply
\(c\sqrt\eta\le\omega(C_D\eta)=o(\sqrt\eta)\), a contradiction.
\end{proof}

The baseline \(\Psi_0\) is itself a translated-source mixture, not
theta. The theorem concerns continuity on that enlarged information
class. It is not a discontinuity statement, a universal upper
Hölder theorem, or an obstruction restricted to the literal theta
source. It leaves open \(d^2\)-scale source resolution, contours
approaching the event with their denominator costs paid, and genuinely
arithmetic signed estimates. The explicit factorial in
\eqref{eq:jet-change} also prevents interpreting fixed-order
control as a free growing-order theorem.

